% Part: many-valued-logic
% Chapter: syntax-and-semantics
% Section: formulas

\documentclass[../../../include/open-logic-section]{subfiles}

\begin{document}

\olfileid[hi]{mvl}{syn}{fml}

\olsection{\usetoken{P}{formula}}

\begin{defn}[सूत्र]
\ollabel{defn:formulas}
किसी प्रतिज्ञप्तिपरक भाषा~$\Lang L$ के \emph{!!{formula}} का समुच्चय
~$\Frm[L]$ आगमनतः इस प्रकार परिभाषित है:
\begin{enumerate}
\item प्रत्येक !!{propositional variable}~$\Obj p_i$ परमाण्विक
  !!{formula} है।
\item ~ $\Lang L$ का प्रत्येक $0$-स्थानी संयोजक (प्रतिज्ञप्तिपरक अचर)
परमाण्विक !!{formula} है।
\item यदि $\star$, ~$\Lang L$ का $n$-स्थानी संयोजक है और $!A_1$, \dots,
$!A_n$ !!{formula} हैं, तो $\star(!A_1, \dots, !A_n)$ !!a{formula} है।
\tagitem{limitClause}{इसके अतिरिक्त कुछ भी !!a{formula} नहीं है।}{}
\end{enumerate}
यदि $\star$ $1$-स्थानी है, तो $\star(!A_1)$ को प्रायः केवल $\star !A_1$
लिखेंगे। यदि $\star$ $2$-स्थानी है, तो $\star(!A_1,!A_2)$ को प्रायः
$(!A_1 \star !A_2)$ लिखेंगे।
\end{defn}

सामान्य रीति से हम प्रायः सबसे बाहरी कोष्ठक बिना बताए छोड़ देंगे।

\begin{ex}
  मानक भाषा~$\Lang{L_0}$ में $\Obj p_1 \lif (\Obj p_1 \land \lnot \Obj
  p_2)$ कोई सूत्र है। गुणनफल तर्कशास्त्र की भाषा में उसे इसके बदले $\Obj
  p_1 \lif (\Obj p_1 \odot \lnot \Obj p_2)$ लिखा जाएगा। यदि भाषा में
  $1$-स्थानी $\triangle$ जोड़ें, तो हमारे पास $\triangle (\Obj p_1 \land
  \Obj p_2) \lif (\triangle \Obj p_1 \land \triangle \Obj p_2)$ जैसे
  सूत्र भी होंगे।
\end{ex}

\end{document}
